% Seventh Philippic --- headnote
% philippic-7, mid-January 43 BC, Rome

\textit{Delivered in the Senate in mid-January 43 BC, while the
legation despatched on 4 January (Sulpicius Rufus, Calpurnius
Piso, Marcius Philippus) was still on the road to Antony at
Mutina, the Seventh Philippic is the holding speech of the
legation interval --- the great anti-peace \textit{sententia}
that prepares the ground for the inevitable rejection of any
terms Antony might bring back. The Senate had been called to
small business: the consul Vibius Pansa on the Appian Way and
the Mint, a tribune of the plebs on the Lupercalia. Cicero uses
those minor referrals as the launching point for the larger
question he insists on keeping alive: while the legation
proceeds, the will to war must not slacken. The mid-January
date is fixed only by month-precision in the manuscript
tradition; nothing in the speech itself permits a sharper
dating.}

\textit{The architecture is the cleanest in the Philippic
series. After the opening preamble (\textsection 1--7), with
its complaint that the legation was sent against Cicero's
judgment, its diagnosis of the rumour-mongers who already
invent Antony's answers and defend them, and its long
encomium of the consul Pansa as the man fitted to the storm
(\textsection 5--7), the speech announces its tripartite
\textit{propositio} at \textsection 8: \textit{ego ille
$\ldots$ pacis semper laudator, semper auctor, pacem cum
M. Antonio esse nolo} --- I, the lifelong author of peace, do
not want peace with Marcus Antonius. The three reasons follow
in order: peace with Antony is \textit{turpis} (shameful,
\textsection 9--15), \textit{periculosa} (dangerous,
\textsection 16--20), and impossible to bring about
(\textit{esse non potest}, \textsection 21--25). Each
demonstration runs in tight rhetorical figure: shamefulness is
shown by the inconsistency of reversing the Senate's many
decrees by which Antony has effectively been judged an enemy
(\textsection 10--13, the great five-fold anaphora of
\textit{non hostem iudicastis}); danger is shown through
the threat of \textit{the Antonii} --- Lucius the murmillo
brother with his catalogue of statue-clientships
(\textsection 16--18, redeploying the set-piece of Philippic 6);
impossibility is shown through the cascade of constituencies
with whom no peace can ever hold --- Senate, Roman equestrians,
the populus Romanus in the contio, the municipalities of
Italy, the named exemplum of Lucius Visidius the Roman
equestrian, Octavian, Decimus Brutus, the province of Gaul
itself (\textsection 21--25). The whole closes with the
peroration of \textsection 26--27: the conditions Antony
would have to meet, with the famous reversal at the end ---
\textit{non illi senatus, sed ille populo Romano bellum
indixerit}, not the Senate but Antony will be declared the
warmaker --- and the warning to Pansa at the tiller of the
ship: do not let so great a preparation come to nothing.}

\textit{The mood is grave but tonally controlled. The
opening \textit{parvis de rebus sed fortasse necessariis
consulimur, patres conscripti} sets the register: the small
business of the day is real, but cannot displace the larger
crisis. The Lucius Antonius set-piece is shorter than in
Philippic 6, but uses the same rhetorical move --- the four
\textit{patron of} anaphora (\textit{patronus} of the
thirty-five tribes, of the equestrian centuries, of the
military tribunes, of the Middle Janus) compresses what
Philippic 6 unfolded statue by statue. The gladiator-language
sharpens (\textsection 17, the Mylasa murmillo who cut down
his Thracian-armed friend in the arena, \textit{luculentam
tamen ipse plagam accepit, ut declarat cicatrix}). The
closing \textit{moriamur} (\textsection 14, \textit{si non
possumus facere $\ldots$ moriamur}) is the rhetorical apex of
the speech and the limit-point of the senatorial register:
let us die, if we cannot wrench these arms away.}

\textit{The speech sits structurally between Philippic 6
(the popular damage-control \textit{contio} of 4 January) and
Philippic 8 (the senatorial \textit{tumultus} speech of
3 February, delivered after the legates' return). Cicero is
holding the line during the interval, preparing the Senate
for what he knows the legation will bring: nothing. Sulpicius
Rufus, the great jurist, died on the road; Piso and Philippus
reached Antony, were rebuffed (Antony added counter-demands:
the consular province of Further Gaul, the veteran lands,
immunity for his acts); by 2 February they were back. On
3 February Cicero delivered Philippic 8 and the Senate at
last decreed \textit{tumultus}. The Seventh Philippic is the
speech by which he held the senatorial nerve in the meantime.}
